\documentclass[11pt,a4paper]{article}
\usepackage[margin=27mm]{geometry}
\usepackage{amsmath,amssymb,amsthm,hyperref}
\newtheorem{lemma}{Lemma}\newtheorem{theorem}[lemma]{Theorem}
\newcommand{\Sz}{\operatorname{Sz}}\newcommand{\F}{\mathbb F}
\title{Extensions of residues containing subfield groups and tori}
\author{Human Compression project}\date{}
\begin{document}\maketitle
\begin{abstract}
We determine all rank-four residually weakly primitive extensions of
the specified residues R10, R11 and R13 in the fixed higher-rank Suzuki
paper. A torus subgroup together with its full Weyl complement determines
its containing subfield groups uniquely. A second argument extracts a
Boolean divisor lattice from residual weak primitivity. Neither argument
assumes the source's classification of all residues.
\end{abstract}

\section{Premises and residual conditions}
Let $S_f=\Sz(2^f)$, with $f$ odd, denote a standard subfield group of
$S_N$, where $f\mid N$. We use the classical subgroup and conjugacy
theorems for Suzuki groups, stated in the accepted preliminary papers.
In particular, every proper nonsoluble subgroup is a simple subfield
group. Simple subfield groups are self-normalizing, and their copies
inside any given Suzuki group form one conjugacy class. The other
maximal subgroups are point stabilizers and normalizers of the three
cyclic torus types. A nonidentity odd-order element determines its
torus uniquely. The three torus orders
\[
 2^f-1,\qquad 2^f+2^{(f+1)/2}+1,\qquad
 2^f-2^{(f+1)/2}+1
\]
are pairwise coprime. Their normalizers have respective complements
of orders $2,4,4$. In each case the complement contains an involution
inverting the torus. Copies of the standard subfield groups intersect
according to $S_a\cap S_b=S_{\gcd(a,b)}$.

The accepted torus construction supplies two compatible nonsplit
families $T_f^+,T_f^-$, together with a common complement of order four.
Their signs need not agree with the signs of the square-root formulas
at every exponent. What matters here is
\[
 T_a^\epsilon\cap T_b^\epsilon=T_{\gcd(a,b)}^\epsilon,
 \qquad
 \gcd(|T_a^\epsilon|,|T_b^\epsilon|)
       =|T_{\gcd(a,b)}^\epsilon|.                 \tag{1}
\]
For $d\mid f$, the two bottom orders divide the two different upper
orders. These are reused structural premises, proved in the accepted
paper \emph{Subfield and torus geometries}.

Write $G_J=\bigcap_{j\in J}G_j$ and $G_\varnothing=G$.
Residual weak primitivity means that for each proper $J\subset I$
some $G_{J\cup\{i\}}$ is maximal proper in $G_J$. It implies
\[
 G_{J\cup\{i\}}<G_J\quad(i\notin J).           \tag{2}
\]
Indeed, equality would imply $G_{I\setminus\{i\}}\le G_i$;
the last panel would then equal $G_I$, contrary to maximal proper
inclusion. The arguments below only need this group condition.

\section{Uniqueness of subfield overgroups}
\begin{lemma}\label{unique}
A fixed simple subfield group is contained in at most one subfield
group of any prescribed degree. The same assertion holds for
$H=C_m:C_c$ when $m>1$ is odd and $C_c$ is the full torus complement:
$c=2$ in a split torus, or $c=4$ in a nonsplit torus.
\end{lemma}
\begin{proof}
The first assertion follows from conjugacy of subfield copies and
self-normalization. More generally, suppose $H\le S_f$ and a conjugate
$gS_fg^{-1}$ also contains $H$. If the copies of $H$ inside $S_f$
are conjugate, adjust $g$ by an element of $S_f$ to normalize $H$.
When $N_G(H)\le S_f$, this forces $gS_fg^{-1}=S_f$.

For the stated torus subgroup, $C_m$ is characteristic: it consists
of the identity and all odd-order elements of $H$. Therefore its
normalizer is contained in the normalizer $T:C_c$ of its unique torus.
Remove the complement component of a normalizing element, leaving
$h\in T$. If $z\in C_c$ is the involution, then
$hzh^{-1}=h^2z$ under a choice of conjugation convention.
Normalization of $H$ implies $h^2\in C_m$. Since $|T|$ is odd,
this implies $h\in C_m$. Hence $N_G(H)=H$.

Copies of $H$ in a Suzuki group are conjugate: the odd subgroup
selects a unique torus type, the cyclic torus has a unique subgroup
of order $m$, and its complements are conjugate in the odd-order
torus extension. This proves the second assertion by the preceding
transporter argument. The use of the full complement is essential;
the assertion does not cover $C_m:C_2$ inside a nonsplit torus.
\end{proof}

\begin{lemma}\label{meet}
Suppose two simple subfield groups of degrees $a,b$ contain the same
$H=C_m:C_c$ of Lemma~\ref{unique}. In the nonsplit case assume $m>5$.
Their intersection is a simple subfield group of degree $d=\gcd(a,b)$.
\end{lemma}
\begin{proof}
The split type persists in every odd-degree subfield containing it:
a nonsplit occurrence would make its odd order divide both
$2^N-1$ and $2^{2N}+1$, whose gcd is one.
In the split case $m\mid\gcd(2^a-1,2^b-1)=2^d-1$.
Thus $H$ has a standard copy in $S_d$. Here $d>1$, since $m>1$.
In the nonsplit case $H$ is nonsplit in both groups, since its
complement has order four. The identity
\[
 \gcd(2^{2a}+1,2^{2b}+1)=2^{2d}+1
\]
holds because $a/d$ and $b/d$ are odd. Thus $m$ divides the product
of the two nonsplit orders at $d$. It divides just one of them:
each bottom order divides a different upper order at $a$, whereas
$m$ divides only the upper order containing its torus. Coprimality
of the upper orders excludes any factor from the other bottom order.
The compatible torus family therefore places a copy of $H$ in $S_d$.
Its action agrees with that of $H$: the order-four action is raising
to $\pm2^d$, and an odd multiple of $d$ gives this action or its
inverse modulo $m$. Reversing a complement generator identifies the
two presentations. Since the only nontrivial nonsplit order at
$d=1$ is $5$, the assumption $m>5$ makes $d>1$.

Conjugate the copy of $H$ in $S_d\le S_a$ onto the given one inside
the first group. In these coordinates the standard $S_b$ is an
overgroup of $H$. Lemma~\ref{unique} identifies it with the second
group. Their intersection is now the standard $S_d$.
\end{proof}

\begin{lemma}\label{field}
Let $G_0=S_{N/c}$ be maximal in $S_N$, where $c$ is prime.
Suppose a proper simple subfield group $K$ shares an embedded
simple subfield with $G_0$, and
$G_0\cap K=S_{N/(ca)}$, where $a$ is prime and the common group
is simple. If $K\nleq G_0$, then $a\ne c$ and $K=S_{N/a}$ after
simultaneous standardization.
\end{lemma}
\begin{proof}
The common simple group and Lemma~\ref{unique} make both groups
standard. If $v$ is the degree of $K$, then
$v\mid N$, $v\nmid N/c$, and
$\gcd(v,N/c)=N/(ca)$. Since $N/c$ lacks just one occurrence of
the prime $c$, the first two conditions force $v_c(v)=v_c(N)$.
Consequently $a\ne c$, as otherwise the gcd would have to lack a
second occurrence of $c$. At $a$ the gcd equality forces
$v_a(v)=v_a(N)-1$; at every other prime different from $c$ it forces
the full valuation in $N$. Thus $v=N/a$. This proof allows $a$ and
$c$ to divide the common base exponent.
\end{proof}

\section{A divisor lattice forced by residual primitivity}
\begin{lemma}\label{divisors}
Let $L=C_M:C_c$ have odd rotation order $M$. Let
$H_i=C_{d_i}:C_c$ for $1\le i\le r$, with a common fixed complement,
and put $\gamma=\gcd(d_1,\ldots,d_r)$. If these intersections satisfy
residual weak primitivity, then there are distinct primes
$\ell_1,\ldots,\ell_r$ such that
\[
 M=\gamma\prod_{i=1}^r\ell_i,
 \qquad d_i=M/\ell_i.                         \tag{3}
\]
Primes $\ell_i$ may divide $\gamma$.
\end{lemma}
\begin{proof}
For $r=1$ the conclusion is precisely maximality; assume $r\ge2$.
Every subgroup between $C_d:C_c$ and $C_M:C_c$ has the form
$C_h:C_c$, $d\mid h\mid M$: it contains the complement, and
removing that component from each element leaves its rotation.
Thus a proper inclusion is maximal exactly when the ratio of
rotation orders is prime. Successively applying residual weak
primitivity gives $r$ prime-index steps from $L$ to the full meet.
Hence $\Omega(M/\gamma)=r$, with multiplicity counted.

For each $i$, (2) says
$\gcd(d_j:j\ne i)>\gamma$. Some prime $\ell_i$ therefore has
\[
 v_{\ell_i}(d_i)<\min_{j\ne i}v_{\ell_i}(d_j).
\]
One prime cannot witness this inequality for two different indices.
We have found $r$ distinct primes dividing $M/\gamma$. Since its
total number of prime factors is $r$, each occurs exactly once and
there are no other prime factors. At the witnessing prime, $d_i$
has the bottom valuation and every other $d_j$ the top valuation.
At all other primes every $d_j$ has the valuation of $M$. These
statements give (3), including when a witnessing prime divides
$\gamma$.
\end{proof}

\section{The torus residues R10 and R11}
The R11 row calls $n,n'$ two primes without expressly distinguishing
them. The prescribed lattice forces them to be distinct: if $n=n'$,
the subgroups $G_{012},G_{023}$ have the same cyclic rotation subgroup
inside $G_{02}$ and contain the same full complement $G_{0123}$.
They coincide, so their meet has order $c\gamma n$ instead of the
prescribed $c\gamma$. Thus the equal-prime case has no realization.
\begin{theorem}\label{extensions}
Let $u=2^e\ge8$ and let $a,b$ be distinct odd primes. Suppose a
rank-four residually weakly primitive pregeometry in $\Sz(q)$ has
a maximal parabolic $G_0=\Sz(u^{ab})$ with the prescribed residue
\[
 G_{01}=\Sz(u^a),\quad G_{03}=\Sz(u^b),\quad
 G_{02}=C_{\gamma nn'}:C_c,
\]
\[
 G_{013}=\Sz(u),\quad G_{012}=C_{\gamma n}:C_c,
 \quad G_{023}=C_{\gamma n'}:C_c,
 \quad G_{0123}=C_\gamma:C_c,                  \tag{4}
\]
where $n,n'$ are distinct primes. For R10 put $c=2$, $\gamma=u-1$;
for R11 put $c=4$ and let $\gamma$ be a nonsplit torus order at $u$.
Then an odd prime $d$ distinct from $a,b$ satisfies
\[
 q=u^{abd},\qquad G_1=\Sz(u^{ad}),\qquad
 G_3=\Sz(u^{bd}),\qquad G_{13}=\Sz(u^d).
\]
There is a prime $n''$ distinct from $n,n'$ such that
\[
 \begin{split}
 G_2&=C_{\gamma nn'n''}:C_c,\\
 G_{12}&=C_{\gamma nn''}:C_c,\qquad
 G_{23}=C_{\gamma n'n''}:C_c,\\
 G_{123}&=C_{\gamma n''}:C_c.
 \end{split}                                \tag{5}
\]
These formulas hold simultaneously after conjugation. In R10,
$n''\mid(u^d-1)/(u-1)$. In R11, $\gamma n''$ divides the compatible
nonsplit torus order at $u^d$, and $\gamma nn'n''$ divides its
ambient order. All remaining intersections are determined by (4)--(5).
\end{theorem}
\begin{proof}
The groups $G_1,G_3$ contain the simple groups $G_{01},G_{03}$,
so are simple subfield groups. All share $G_{013}$, allowing
simultaneous standardization by Lemma~\ref{unique}. Maximality
of $G_0$ gives $q=(u^{ab})^d$ for a prime $d$. Applying
Lemma~\ref{field} to its two specified maximal subfield
intersections gives $d\ne a,b$, and the asserted degrees and meet.
Conjugate the common full torus normalizer $G_{0123}$ inside $S_e$
onto the chosen standard normalizer. This preserves all three
standard simple groups and fixes their torus coordinates simultaneously.

Put $H=G_{02}$. If $G_2$ were a simple subfield group,
Lemma~\ref{meet} would make $G_0\cap G_2$ a simple subfield
group, contradicting $G_0\cap G_2=H$. In the nonsplit case its
rotation order $\gamma nn'$ exceeds five even when $\gamma=5$.
This explains why one must use the entire $H$, rather than claim
that its base rotation $C_\gamma$ generates the base field.

Choose a minimal Suzuki overgroup of $G_2$, allowing the ambient
group. Since $G_2$ is not itself a Suzuki group, it lies in a
maximal subgroup of this overgroup. Minimality rules out the
subfield type, so it lies in a point stabilizer or torus normalizer.
A point stabilizer cannot contain $H$: its quotient by its normal
two-subgroup is cyclic of odd order, whereas the involution in $H$
inverts a nontrivial odd rotation. The images of that involution
and its inverse action would force this rotation to equal its
inverse. It cannot disappear in the two-subgroup either.
The wrong torus types are excluded by coprimality. The unique
torus containing the odd part of $H$ is therefore also the torus
for $G_2$. Thus $G_2=C_M:C_c$, with the same complement as $H$.

Its three parabolic intersections with $G_0,G_1,G_3$ all contain
the complement from $G_{0123}$. Their full meet is
$C_\gamma:C_c$. Lemma~\ref{divisors} says that $M/\gamma$ is the
product of three distinct primes, one omitted at each parabolic.
The prescribed orders of $G_{02},G_{012},G_{023}$ identify two
as $n,n'$ and force the third, $n''$, in precisely the positions
of (5). Intersecting with $G_{13}=\Sz(u^d)$ proves the final
divisibility condition. The inherited unique torus determines its
compatible branch. Its containment in the ambient torus also
gives the asserted ambient divisibility.
\end{proof}

\section{The pure subfield residue R13}
\begin{theorem}
Suppose $G_0=S_{N/d}$ is maximal in $S_N$ and its prescribed rank-three
residue has parabolics $G_{0i}=S_{N/(da_i)}$ for $i=1,2,3$,
where $a_1,a_2,a_3$ are distinct odd primes and their common
intersection is simple. Every rank-four residually weakly primitive
extension has $a_i\ne d$ and $G_i=S_{N/a_i}$, after one simultaneous
conjugation. All intersections are the standard subfield intersections.
\end{theorem}
\begin{proof}
Each $G_i$ contains the simple $G_{0i}$, hence is a simple subfield
group. The common simple subgroup standardizes every one of them
by Lemma~\ref{unique}. Strictness excludes $G_i\le G_0$.
Lemma~\ref{field} now gives the stated degrees. Their simultaneous
standard coordinates give every intersection by the gcd formula.
\end{proof}

\section{Source scope and attribution}
The source is D.~Leemans, \emph{The residually weakly primitive
pre-geometries of the Suzuki simple groups}, Note di Matematica
20(1) (2000/2001), 1--20, fixed edition, pp.~14--15.
The R10 paragraph interchanges the labels of its second new Suzuki
parabolic and its dihedral parabolic. The consistent labels are
$G_1,G_3$ for the Suzuki groups and $G_2$ for the dihedral group,
as fixed by the prescribed intersections in its opening sentence.

The ambient structure and torus compatibility are borrowed as stated.
The overgroup, divisor-lattice and conditional extension arguments
above are supplied here to repair a gap in our first claim ledger.
A refutation of the source's global classification did not settle
these distinct conditional assertions. The present theorems settle
the rank-four assertions for the actual prescribed residues. The
accepted constructions already prove realization of their admissible
canonical parameters. Neither those existence results nor this note
classifies all unrestricted higher-rank relative positions. In
particular, the source's separate upper rank bound for pure subfield
constructions remains refuted by the previously accepted repeated-prime
example. No new open-conjecture, principal-refutation, Lean or
compression credit is asserted here.
\end{document}
